\section{Encoding sparse polynomial systems}\label{sec:tensor-encoding}
We prove Lemma~\ref{lem:tensor-sparse-query}. The first step represents
each equation by an integer polynomial modulo a power of two. Truncating
the product across all equations keeps its total input degree linear
in the number of variables. The proof in
Section~\ref{subsec:tensor-sparse-proof} then encodes its monomials
as tensor indices and applies Theorem~\ref{thm:tensor-contraction}.

Recall the common-zero count $N(\mathcal P)$ from
\eqref{eq:common-zero-count}.
In the lift below, $m_{\mathrm{test}}$ denotes the number of equations,
and each $P_i$ has degree at most $k$. The list of its monomials
is supplied over $\mathbb F_2$.

\subsection{An integer polynomial for one equation}\label{subsec:rational-lift}
On Boolean inputs we may replace each positive power $x_v^j$ by $x_v$.
Combine duplicate monomials modulo two. If $\mathcal M_i$ is the
resulting list of variable sets for $P_i$, define its integer lift by
\[
 L_i(\vec X)=\sum_{M\in\mathcal M_i}\prod_{v\in M}X_v .
\]
The empty monomial is $1$ and the empty sum is $0$. The value $L_i(x)$
counts the monomials equal to one, so $P_i(x)=0$ precisely when
$L_i(x)$ is even. Write $I_i=\bigcup_{M\in\mathcal M_i}M$. Its scope size is at most $k$ times its monomial-list length,
while its degree is at most $k$.

Set $\ell=n+1$. Since $0\le N(\mathcal P)\le2^n<2^\ell$, it suffices
to obtain its residue modulo $2^\ell$. We first represent the
evenness indicator by an integer polynomial modulo this power of two.
Multiplying $m_{\mathrm{test}}$ such indicators directly would give
only a quadratic degree bound; the form of the next lemma will let us
control the degree across all equations together.

Put $h(Y)=Y(Y-1)$. Its value is even at every integer $Y$.
Writing the indicator in powers of $h$ exposes one factor of two for
each copy of $h$ in the product.

\begin{lemma}\label{lem:rational-idempotent}
For every integer $\ell\ge1$, one can compute polynomials
$e_0,\ldots,e_{\ell-1}\in\mathbb Z[Y]$, each of degree at most one,
such that
\begin{equation}\label{eq:idempotent-expansion}
 p_\ell(Y)=\sum_{j=0}^{\ell-1}e_j(Y)h(Y)^j,
 \qquad
 p_\ell(y)\equiv\mathds{1}_{\{2\mid y\}}\pmod{2^\ell}
 \quad(y\in\mathbb Z).
\end{equation}
The coefficients of the $e_j$ have $\poly(\ell)$ bits, and the
construction uses $\poly(\ell)$ bit operations and $\poly(\ell)$
bits of space.
\end{lemma}
\begin{proof}
Let $F_\ell$ be the modulus-amplifying polynomial of Beigel and
Tarui~\cite[Section~2.3]{BT94}. Chan and Williams apply $F_\ell$ to
sums of clause-violation indicators to count satisfying assignments
deterministically~\cite[Theorem~2.5 and proof of Theorem~1.6]{CW21}.
We need the indicator of even parity, so we take its complement:
\begin{equation}\label{eq:binary-indicator}
 p_\ell(Y)=1-F_\ell(Y)=(1-Y)^\ell
       \sum_{j=0}^{\ell-1}\binom{\ell+j-1}{j}Y^j.
\end{equation}
The finite sum is the inverse of $(1-Y)^\ell$ modulo $Y^\ell$.
Indeed, the coefficient of $Y^j$ in $(1+Y+Y^2+\cdots)^\ell$
counts $\ell$-tuples of nonnegative integers summing to $j$ and equals
$\binom{\ell+j-1}{j}$. Therefore
\[
 p_\ell(Y)\equiv1\pmod{Y^\ell},
 \qquad
 p_\ell(Y)\equiv0\pmod{(1-Y)^\ell}.
\]
For even $y$, the first congruence gives $p_\ell(y)\equiv1\pmod{2^\ell}$.
For odd $y$, the second gives $p_\ell(y)\equiv0\pmod{2^\ell}$.

The polynomial $p_\ell$ has integer coefficients and degree less
than $2\ell$. Successive division by the monic quadratic
$h=Y^2-Y$ writes it uniquely as
$\sum_{j<\ell}e_jh^j$, with integer remainders $e_j$ of degree
less than two. Lemma~\ref{lem:rational-lift-height} in
Appendix~\ref{app:arithmetic} bounds their coefficients and proves
the $\poly(\ell)$ bit-operation bound for this construction.
\end{proof}

\subsection{Keeping the product degree linear}\label{subsec:total-order}
The full product $\prod_{i=1}^{m_{\mathrm{test}}} p_\ell(L_i(\vec{X}))$ has degree at most
$km_{\mathrm{test}}(2\ell-1)$. Each factor $h(y)$ contributes a factor of two at
integer arguments, so a term containing at least $\ell$ such factors
can be discarded modulo $2^\ell$. This lets us retain only terms
whose total power of $h$ is less than $\ell$, across all equations together.
Introduce a formal variable $V$ to record that power, using the same
$V$ for every equation. Its exponent records a guaranteed power of two
dividing the corresponding term at integer arguments. It is not an
additional Boolean variable to be summed over.
For the rational computation we use the coefficient algebra
\[
 \mathcal E_\ell=\mathbb Q[V]/(V^\ell).
\]
In this algebra every multiplication discards powers $V^s$ with
$s\ge\ell$. Once the product is formed, we add its retained
coefficients using the linear map
\[
 \mathcal S_\ell\!\left(\sum_{s=0}^{\ell-1}b_sV^s\right)
 =\sum_{s=0}^{\ell-1}b_s .
\]
This map need not preserve products, so it is applied only after the
factors have been multiplied. We use the same notation for polynomial
coefficients $b_s$.

\begin{lemma}\label{lem:total-order}
Let $m_{\mathrm{test}}\ge0$ and $\ell\ge1$ be integers. For integers
$y_1,\ldots,y_{m_{\mathrm{test}}}$, put
\[
 B_i(V)=\sum_{j=0}^{\ell-1}e_j(y_i)h(y_i)^jV^j.
\]
Then
\[
 \mathcal S_\ell\!\Bigl(\prod_{i=1}^{m_{\mathrm{test}}} B_i(V)\bmod V^\ell\Bigr)
 \equiv\prod_{i=1}^{m_{\mathrm{test}}}\mathds{1}_{\{2\mid y_i\}}\pmod{2^\ell}.
\]
Independently, for any nonnegative integer $d_*$, if the $y_i$ are
polynomials over $\mathbb Q$ in any common collection of inputs,
each of total degree at most $d_*$,
then each retained coefficient of $V^s$, and hence their sum, has
total input degree at most
$d_*(m_{\mathrm{test}}+2(\ell-1))$.
\end{lemma}
\begin{proof}
Before truncation, a term contributing to the coefficient of $V^s$
has the form
\[
 \prod_{i=1}^{m_{\mathrm{test}}} e_{j_i}(y_i)h(y_i)^{j_i},
 \qquad \sum_{i=1}^{m_{\mathrm{test}}} j_i=s.
\]
At integer arguments it is divisible by $2^s$. Dropping every
coefficient with $s\ge\ell$ therefore leaves the sum of coefficients
unchanged modulo $2^\ell$. Without truncation, that sum is
$\prod_i p_\ell(y_i)$, which has the required congruence by
Lemma~\ref{lem:rational-idempotent}.
Truncating after each multiplication gives the same result as
truncating the full product: every later factor has nonnegative
$V$-degree, so a discarded term can never return to degree below
$\ell$. This also explains why the coefficient sum is taken at
the end; it is the retained total orders that have the required
congruence.

For polynomial arguments, every retained term has
$\sum_i j_i\le\ell-1$. Since $\deg e_{j_i}\le1$ and $\deg h=2$,
its degree is at most
\[
 \sum_{i=1}^{m_{\mathrm{test}}} d_*(1+2j_i)
 \le d_*(m_{\mathrm{test}}+2(\ell-1)).
\]
The bound holds coefficient by coefficient, as well as for their sum.
\end{proof}

To see why the truncation uses the total order across all equations,
take $\ell=2$. Then
\[
 p_2(Y)=1-3Y^2+2Y^3=(1-Y)+(2Y-1)h(Y).
\]
For two equations, multiplication modulo $V^2$, followed by
$\mathcal S_2$, retains
\[
 (1-y_1)(1-y_2)
 +(2y_1-1)h(y_1)(1-y_2)
 +(1-y_1)(2y_2-1)h(y_2).
\]
Its degree is at most four. The omitted term
$(2y_1-1)(2y_2-1)h(y_1)h(y_2)$ has degree six, but is divisible
by four at integer inputs.


\subsection{From the retained monomials to a network}\label{subsec:tensor-sparse-proof}
\begin{proof}[Proof of Lemma~\ref{lem:tensor-sparse-query}]
The network will sum coefficients of incidence monomials, with
test vertices supplying those coefficients and variable vertices
supplying their sums over Boolean values.
Put $\lambda=\lceil\log_2 n\rceil$ and
$\nu=\lceil\log_2\lambda\rceil$.
We will apply Theorem~\ref{thm:tensor-contraction} with $s=n\nu$,
so we need $O_{k,c}(n\nu)$ wires and maximum degree at most
$\lceil(n\nu)^{9/10}\rceil$. Each vertex will stand for a
group of tests or of free variables, and its degree will be at most the
number of incidences in its group. Each test already reads only
polylogarithmically many variables, but a single variable may belong to
all $a$ scopes. We first fix the high-incidence variables in all
possible ways. Afterwards every free variable lies in only
polylogarithmically many scopes. After grouping, this will bound the
degrees on both sides by $n^{3/4}\lambda^{O_{k,c}(1)}$, below the
required maximum for all sufficiently large $n$.
Within each monomial replace positive powers of a variable by its
first power, and combine duplicate monomials modulo two. The scope
of test $i$ consists of the variables in its remaining monomials.
It has size at most $k\lambda^c$.
Let $N_0$ be the sum of the scope sizes, counting a pair $(i,v)$
once if variable $v$ belongs to test $i$. Thus $N_0\le ak\lambda^c\le ckn\lambda^c$, since $a\le cn$.
Fix, in all possible ways, the variables belonging to more than
\[
 D_*=\max\{1,\lceil N_0/n\rceil\}\lceil\sqrt{\lambda}\rceil
\]
scopes. There are at most $n/\sqrt{\lambda}$ such variables. Their
assignments partition the original count into at most
$2^{n/\sqrt{\lambda}}$ branch counts. On each branch substitute the
fixed bits and combine duplicate monomials again. The free scope
$I_i$ of test $i$ still has size at most $k\lambda^c$, and each free
variable appears in at most $D_*=\lambda^{O_{k,c}(1)}$ scopes. Free
variables absent from every scope remain declared. A constant-one
equation makes the branch count zero. We may discard constant-zero
equations, and if none remain the answer for the branch is two to the
power of its number of free variables. Below $a$ denotes the number
of remaining equations and $n$ the original variable count.

\paragraph{The polynomial on the incidences.}
Give each pair $(i,v)$ with $v\in I_i$ its own indeterminate
$z_{iv}$. Write $z_i=(z_{iv})_{v\in I_i}$. Form
$\widetilde L_i(z_i)$ by summing the remaining monomials of the
normalized branch polynomial with integer coefficients, using $z_{iv}$
for $x_v$. We use the ring
\[
 \mathcal B=(\mathbb Z/2^p\mathbb Z)[V]/(V^\ell),
 \qquad p=\ell=n+1 .
\]
Take $h(Y)=Y(Y-1)$ and the integer polynomials $e_j$ from
Lemma~\ref{lem:rational-idempotent}. Reduce their coefficients
modulo $2^p$, and form the factor
\[
 B_i(z_i,V)=\sum_{j=0}^{\ell-1}
   e_j(\widetilde L_i(z_i))\,h(\widetilde L_i(z_i))^j V^j .
\]
Let $\overline B_i$ be its multilinear representative in the $z_i$,
obtained in each coefficient of $V$ by replacing each positive
power of an incidence indeterminate by its first power. It agrees
with $B_i$ on every Boolean assignment to $z_i$. Define
$F(z,V)=\prod_i\overline B_i(z_i,V)$ in $\mathcal B[z]$. The
incidence indeterminates for different tests are disjoint, so $F$
is multilinear even before substituting the same Boolean $x_v$ for
all copies of variable $v$.

The coefficient of $V^j$ in $\overline B_i$ has incidence degree
at most $k(1+2j)$, since $\deg e_j\le1$, $\deg h=2$, and
multilinearization only lowers monomial degrees. Multiplication in
$\mathcal B$ retains only terms with $\sum_i j_i\le\ell-1$.
Consequently every incidence monomial in $F$ has degree at most
\begin{equation}\label{eq:sharp-incidence-degree}
 K=k\bigl(a+2(\ell-1)\bigr)=O_{k,c}(n).
\end{equation}
Over $\mathcal B$, use the same coefficient-sum formula
$\mathcal S_\ell(\sum_{j<\ell}b_jV^j)=\sum_{j<\ell}b_j$, now with
$b_j\in\mathbb Z/2^p\mathbb Z$. At $z_{iv}=x_v$, the multilinearized factor
$\overline B_i$ agrees with $B_i$ for the normalized branch polynomial.
The integer $\widetilde L_i((x_v)_{v\in I_i})$ has the same parity
as the original $P_i$ evaluated at $x$ and the fixed bits. Applying
Lemma~\ref{lem:total-order} to these integers shows that the branch
count is congruent modulo $2^p$ to
\begin{equation}\label{eq:sharp-branch-residue}
 \mathcal S_\ell\left(\sum_x F\bigl((z_{iv}=x_v)_{iv},V\bigr)\right),
\end{equation}
where $x$ runs over all free Boolean assignments. We apply the
coefficient-sum map after multiplying the test factors in the
truncated ring.

\paragraph{Indexing the monomials.}
Giving each incidence its own wire would use up to
$N_0\le ckn\lambda^c$ wires, while the claimed exponent allows only
$O_{k,c}(n\nu)$. A monomial of $F$ selects at most $K$ incidences.
We group the tests and the free variables, fix how many selected
incidences join each pair of groups, and let the wires between two
groups index the subsets of that size. The sizes are enumerated
separately as an array $(d_{gh})$, with one entry per pair of groups;
by \eqref{eq:sharp-wire-budget} below the binary subset indices give
the $O_{k,c}(n\nu)$ wire bound. We choose
$J=\lceil n^{1/4}\rceil$ groups on each side so that the $J^2$
array entries cost only $2^{O_{k,c}(\sqrt n\log n)}$ choices,
while each group has few enough incidences to meet the
maximum-degree bound.
Partition the tests and the free variables separately into these
$J$ groups, with sizes differing by at most one on each side. Write $\mathcal T_g$ for a test group and
$\mathcal X_h$ for a variable group. A group contains
$O_c(n^{3/4})$ objects, so the number of its incidences is at most
$n^{3/4}\lambda^{O_{k,c}(1)}$ on either side. Let $N_{gh}$ be the
number of incidences from $\mathcal T_g$ to $\mathcal X_h$ and let
$N=\sum_{g,h}N_{gh}\le N_0$.

For an incidence monomial let $d_{gh}$ be the number of its
selected incidences joining groups $g,h$. Enumerate arrays of
integers $(d_{gh})$ with $0\le d_{gh}\le N_{gh}$ and
$\sum_{g,h}d_{gh}\le K$. There are at most
\[
 (K+1)^{J^2}=2^{O_{k,c}(\sqrt n\log n)}
\]
arrays, which can be enumerated using $J^2$ counters. For a fixed
array there are
\[
 C_{gh}=\binom{N_{gh}}{d_{gh}}
\]
subsets of the specified size for this pair. Draw
$\lceil\log_2 C_{gh}\rceil$ Boolean wires between the two group
vertices to index them in lexicographic order. For $C_{gh}=1$
both endpoints use the unique subset (possibly the entire incidence
set) with no wires. Any unused bit string gives tensor entry zero.
To decode a valid index, take the incidences in order and count the
possible completions with or without each incidence using binomial
coefficients. These have polynomial bit length, so decoding uses
polynomial time and space.

The union of the subsets over all pairs specifies a squarefree
monomial of degree at most $K$ in $N$ indeterminates. Including
all monomials of this degree bound gives the upper bound
$\binom{N+K}{K}$ by stars and bars. Thus the number
of wires $M$ in the network for this array satisfies
\begin{align}
 M &\le J^2+\log_2\prod_{g,h} C_{gh}
      \le J^2+\log_2\binom{N+K}{K} \notag\\
   &\le J^2+K\log_2\!\left(\frac{\mathrm e(N+K)}{K}\right)
      = O_{k,c}(n\nu).
 \label{eq:sharp-wire-budget}
\end{align}
Here $K=\Theta_{k,c}(n)$ and $N\le n\lambda^{O_{k,c}(1)}$.
Also $\lceil\log_2 C_{gh}\rceil\le N_{gh}$, so the degree of each
vertex is at most its number of incidences,
$n^{3/4}\lambda^{O_{k,c}(1)}$.

\paragraph{Tensor entries and contraction.}
A valid assignment of wire bits determines a subset $S_i\subseteq I_i$
at each test. Write $z_{i,S_i}=\prod_{v\in S_i}z_{iv}$. The entry
at a test group $g$ is the coefficient of its indicated monomial
in $\prod_{i\in\mathcal T_g}\overline B_i(z_i,V)$. Since the $z_i$
are disjoint, this is
\[
 \prod_{i\in\mathcal T_g}[z_{i,S_i}]\overline B_i(z_i,V)
 \qquad\text{in }\mathcal B.
\]
For a variable group $h$, put $b_v=|\{i:v\in S_i\}|$
for $v\in\mathcal X_h$. Its entry is
\[
 \sum_{x\in\{0,1\}^{\mathcal X_h}}
        \prod_{v\in\mathcal X_h}x_v^{b_v}
   =2^{|\{v\in\mathcal X_h:b_v=0\}|}\pmod{2^p}.
\]
Indeed a positive exponent requires $x_v=1$ and a zero exponent
leaves $x_v$ free. This includes variables outside all test scopes.

Each monomial of $F$ occurs in exactly one array and has exactly
one valid index at each group pair. Its coefficient in $F$ is the
product of the test-group coefficients. After the incidences with
the same $v$ are set equal, its sum over Boolean assignments is
the product of the variable-group entries. By
\eqref{eq:sharp-incidence-degree}, summing over these arrays exhausts
all monomials. Consequently the element of $\mathcal B$ inside
$\mathcal S_\ell$ in \eqref{eq:sharp-branch-residue} is the sum of
the contractions of these bipartite networks.

To evaluate one test-side coefficient, let $\mathbf1_T$ be the
Boolean vector on $I_i$ with support $T\subseteq S_i$. At
$\mathbf1_T$, the multilinear polynomial evaluates to the sum of its
coefficients indexed by subsets of $T$, and it agrees there with
$B_i$. Inverting this subset sum gives
\[
 [z_{i,S_i}]\overline B_i(z_i,V)=
 \sum_{T\subseteq S_i}(-1)^{|S_i|-|T|}B_i(\mathbf1_T,V)
 \qquad\text{in }\mathcal B.
\]
Each term costs polynomially many bit operations. Evaluate
$\widetilde L_i(\mathbf1_T)$ from the listed monomials and use the
polynomials $e_j$ of Lemma~\ref{lem:rational-idempotent}, which
are constructible in polynomial time. Store the $\ell$ coefficients
of $B_i(\mathbf1_T,V)$ modulo $2^p$. There are at most
$2^{|S_i|}\le 2^{k\lambda^c}$ subsets to process. Sum one test
at a time and multiply the resulting coefficients in $\mathcal B$
over the group. This evaluates a test-group entry in
$2^{\lambda^{O_{k,c}(1)}}$ bit operations and polynomial space.
The variable-side entries and checks for invalid wire indices take
polynomial time. The graph, degree array, and entry procedures have
polynomial-size descriptions.

Apply Theorem~\ref{thm:tensor-contraction} with size parameter
$s=n\nu$, root precision $p=n+1$, and truncation order $\ell=n+1$.
The graph has $2J$ vertices and $O_{k,c}(s)$ wires by
\eqref{eq:sharp-wire-budget}. For all sufficiently large $n$ its
maximum degree is at most $\lceil s^{9/10}\rceil$, its precision
and truncation order are $O(s)$, and its entry-evaluation time is
at most $2^{O_{k,c}(s^{9/10})}$. Each contraction then takes
$n^{O_{k,c}(1)}2^{O_{k,c}(s/\sqrt{\log s})}$ bit operations and
polynomial space. Here $s/\sqrt{\log s}=O(n\nu/\sqrt\lambda)$.
The high-incidence branches and degree arrays multiply the
running time by at most
\[
 2^{\,n/\sqrt{\lambda}+O_{k,c}(\sqrt n\log n)}.
\]
Both terms in this exponent are $O_{k,c}(n\nu/\sqrt{\lambda})$.
Process the branches and arrays sequentially. On each branch,
sum the contractions in $\mathcal B$ and apply $\mathcal S_\ell$.
Adding the branch residues modulo $2^{n+1}$ gives the original
count, since it lies between $0$ and $2^n$. This proves the
claimed bounds, using direct enumeration for bounded exceptional
sizes.
\end{proof}

